% GENERATED FILE. Edit canonical lessons, then run npm run generate:books.
% canonical-source-hash: fnv1a64:479f9c1b6d4d3e2f
% canonical-lessons: ZH-C397-fanggezi, ZH-C397-nandao, ZH-W397-nan-hard, ZH-C397-nande, ZH-C397-weinan, ZH-C397-jiekou

\chapter{Surely Not}
\label{ch:zh-surely-not}
\hlchaptermodality{397}

\begin{chapteropening}
\textbf{By the end of this chapter:} \emph{I can talk about being stood up, asking whether something can really be so, rare chances, awkward spots and excuses.}

Complete the last lesson of chapter 397: I can talk about being stood up, asking whether something can really be so, rare chances, awkward spots and excuses.
\end{chapteropening}

\section[fàng gēzi]{fàng gēzi — to stand someone up}

\label{lesson:ZH-C397-fanggezi}



\begin{quote}
Before the new one: say the Mandarin for round the clock, then the Mandarin for even if.
\end{quote}

\subsection*{You'll want to know: fàng gēzi}
\emph{fàng gēzi} — \textquotedblleft{}to stand someone up\textquotedblright{}.

\begin{quote}
\emph{rúguǒ tā yòu fàng gēzi wǒ jiù bù děng tā le} — If he stands me up again, I won't wait for him any more.
\end{quote}

Its characters come later; for now it is a word for the ear.

\begin{cousinweb}
\emph{fàng gēzi} is built from \emph{fàng} 'to let go' and \emph{gēzi} 'a pigeon', all inherited from Middle Chinese, \emph{pjangH}, \emph{kop} and \emph{tsiX} (Guangyun, in Baxter's transcription), which glosses the second 'the name of a bird'. Letting the pigeon fly is a modern saying, and how it came to mean standing someone up is not settled.
\end{cousinweb}

\begin{grammarlens}[title={What you've built}]
One more word for saying what would happen.
\end{grammarlens}

\subsection*{Your turn}
\begin{itemize}
\raggedright
  \item \emph{Say it:} \emph{fàng gēzi}
  \item \emph{Say it:} \emph{fàng gēzi}, once more
  \item \emph{Say it:} \emph{rúguǒ tā yòu fàng gēzi wǒ jiù bù děng tā le} — If he stands me up again, I won't wait for him any more.
  \item \emph{Recall:} say the Mandarin for round the clock, then the Mandarin for even if, then say \emph{fàng gēzi} again
\end{itemize}

\begin{tcolorbox}[breakable,colback=teal!4,colframe=teal!35!black,title={Before you move on}]
What does fàng gēzi mean? (\textquotedblleft{}To stand someone up\textquotedblright{}.) Say it once more.
\end{tcolorbox}
\section[\texorpdfstring{\zh{难道} (nándào)}{难道 (nándào)}]{\zh{难道} (nándào) — surely not (in a question)}

\label{lesson:ZH-C397-nandao}



\begin{quote}
Before the new one: say the Mandarin for even if, then the Mandarin for to stand someone up.
\end{quote}

\subsection*{You'll want to know: \zh{难道}}
\textbf{\zh{难道}} — \emph{nándào} — \textquotedblleft{}surely not (in a question)\textquotedblright{}.

\begin{quote}
\textbf{\zh{难道你不想去吗}} \emph{nándào nǐ bù xiǎng qù ma} — Surely you want to go?
\end{quote}

The new character, in pieces: \zh{难} — ten strokes in two pieces: on the left \zh{又} made narrow, with a dot for its last stroke; on the right the bird piece, eight strokes. The next lesson puts it in your hand.

\begin{cousinweb}
\textbf{\zh{难道}} (\emph{nándào}) is built from \zh{难} 'hard' and \emph{dào} 'to say', both inherited from Middle Chinese, \emph{nan} and \emph{dawX} (Guangyun, in Baxter's transcription), which glosses the first 'difficult' and the second 'a way'. The sense 'to say' is not given there, and how the pair came to ask whether something can really be so is not settled.
\end{cousinweb}

\begin{grammarlens}[title={What you've built}]
One more word for saying what would happen.
\end{grammarlens}

\subsection*{Your turn}
\begin{itemize}
\raggedright
  \item \emph{Say it:} \emph{nándào}
  \item \emph{Say it:} \emph{nándào}, once more
  \item \emph{Say it:} \emph{nándào nǐ bù xiǎng qù ma} — Surely you want to go?
  \item \emph{Recall:} say the Mandarin for even if, then the Mandarin for to stand someone up, then say \emph{nándào} again
\end{itemize}

\begin{tcolorbox}[breakable,colback=teal!4,colframe=teal!35!black,title={Before you move on}]
What does \zh{难道} mean? (\textquotedblleft{}Surely not (in a question)\textquotedblright{}.) Say it once more.
\end{tcolorbox}
\section[\texorpdfstring{\zh{难} (nán)}{难 (nán)}]{\zh{难} — a narrow \zh{又}, then the bird piece}

\label{lesson:ZH-W397-nan-hard}



\begin{quote}
The word first, then the character inside it.

\begin{itemize}
\raggedright
  \item \emph{Recall:} say \emph{nándào} — \textbf{R1}, one lesson back, before it decays
\end{itemize}

\begin{quote}
\zh{取} \zh{又}
\end{quote}

\emph{Read it:} each one aloud, with nothing said first — \emph{qǔ}, \emph{yòu}
\end{quote}

\subsection*{Script — a narrow \zh{又}, then the bird piece}
\hlblockfigure{\detokenize{figures/ZH-W397-nan-hard-filmstrip.pdf}}{How \zh{难} is written, stroke by stroke}

\begin{quote}
\zh{难}
\end{quote}

\textbf{Two pieces, left before right.}

\textbf{The left piece, \zh{又} made narrow, two strokes:} a bar that turns and falls to the lower left, then a dot instead of the long stroke.

\textbf{The right piece, the bird piece, eight strokes.} A stroke falling to the left. A long vertical down from it. A dot at the top right. Three short bars, one under another. A second vertical down through them. Then a long bar across the foot.

2 + 8 = ten.

\subsection*{Your turn}
\begin{itemize}
\raggedright
  \item \emph{Write it:} \zh{难}, all ten, with the model in front of you
  \item \emph{Cover:} write it again from memory, and count the pieces, not the strokes
\end{itemize}

\begin{tcolorbox}[breakable,colback=teal!4,colframe=teal!35!black,title={Before you move on}]
How many strokes in \zh{难}? (\textbf{Ten}.) What ends the left piece? (\textbf{A dot}.) How many short bars in the bird piece before the long one? (\textbf{Three}.)

Source: \href{https://raw.githubusercontent.com/chanind/hanzi-writer-data/68d10a4b21150cae5e1ebbd223eed289cf32d90c/data/\%E9\%9A\%BE.json}{Hanzi Writer Data \zh{难}, pinned stroke paths and medians}.
\end{tcolorbox}
\section[\texorpdfstring{\zh{难得} (nándé)}{难得 (nándé)}]{\zh{难得} (nándé) — rare; hard to come by}

\label{lesson:ZH-C397-nande}



\begin{quote}
Before the new one: say the Mandarin for to stand someone up, then the Mandarin for surely not (in a question).
\end{quote}

\subsection*{You'll want to know: \zh{难得}}
\textbf{\zh{难得}} — \emph{nándé} — \textquotedblleft{}rare; hard to come by\textquotedblright{}.

\begin{quote}
\textbf{\zh{我们难得有时间一起吃饭}} \emph{wǒmen nándé yǒu shíjiān yìqǐ chī fàn} — We seldom get time to eat together.
\end{quote}

\begin{cousinweb}
\textbf{\zh{难得}} (\emph{nándé}) is built from \zh{难} 'hard' and \zh{得} 'to get', both inherited from Middle Chinese, \emph{nan} and \emph{tok} (Guangyun, in Baxter's transcription), which glosses the first 'difficult'.
\end{cousinweb}

\begin{grammarlens}[title={What you've built}]
One more word for saying what would happen.
\end{grammarlens}

\subsection*{Your turn}
\begin{itemize}
\raggedright
  \item \emph{Say it:} \emph{nándé}
  \item \emph{Say it:} \emph{nándé}, once more
  \item \emph{Say it:} \emph{wǒmen nándé yǒu shíjiān yìqǐ chī fàn} — We seldom get time to eat together.
  \item \emph{Recall:} say the Mandarin for to stand someone up, then the Mandarin for surely not (in a question), then say \emph{nándé} again
\end{itemize}

\begin{tcolorbox}[breakable,colback=teal!4,colframe=teal!35!black,title={Before you move on}]
What does \zh{难得} mean? (\textquotedblleft{}Rare; hard to come by\textquotedblright{}.) Say it once more.
\end{tcolorbox}
\section[\texorpdfstring{\zh{为难} (wéinán)}{为难 (wéinán)}]{\zh{为难} (wéinán) — in an awkward spot; torn (about what to do)}

\label{lesson:ZH-C397-weinan}



\begin{quote}
Before the new one: say the Mandarin for surely not (in a question), then the Mandarin for hard to come by.
\end{quote}

\subsection*{You'll want to know: \zh{为难}}
\textbf{\zh{为难}} — \emph{wéinán} — \textquotedblleft{}in an awkward spot; torn (about what to do)\textquotedblright{}.

\begin{quote}
\textbf{\zh{我不想说不所以很为难}} \emph{wǒ bù xiǎng shuō bù suǒyǐ hěn wéinán} — I don't want to say no, so I'm in an awkward spot.
\end{quote}

\begin{cousinweb}
\textbf{\zh{为难}} (\emph{wéinán}) is built from \zh{为} 'to make' and \zh{难} 'hard', both inherited from Middle Chinese, \emph{hjwe} and \emph{nan} (Guangyun, in Baxter's transcription), which glosses the first 'to do' and the second 'difficult'.
\end{cousinweb}

\begin{grammarlens}[title={What you've built}]
One more word for saying what would happen.
\end{grammarlens}

\subsection*{Your turn}
\begin{itemize}
\raggedright
  \item \emph{Say it:} \emph{wéinán}
  \item \emph{Say it:} \emph{wéinán}, once more
  \item \emph{Say it:} \emph{wǒ bù xiǎng shuō bù suǒyǐ hěn wéinán} — I don't want to say no, so I'm in an awkward spot.
  \item \emph{Recall:} say the Mandarin for surely not (in a question), then the Mandarin for hard to come by, then say \emph{wéinán} again
\end{itemize}

\begin{tcolorbox}[breakable,colback=teal!4,colframe=teal!35!black,title={Before you move on}]
What does \zh{为难} mean? (\textquotedblleft{}In an awkward spot; torn (about what to do)\textquotedblright{}.) Say it once more.
\end{tcolorbox}
\section[jièkǒu]{jièkǒu — excuse (a poor reason)}

\label{lesson:ZH-C397-jiekou}



\begin{quote}
Before the new one: say the Mandarin for hard to come by, then the Mandarin for in an awkward spot.
\end{quote}

\subsection*{You'll want to know: jièkǒu}
\emph{jièkǒu} — \textquotedblleft{}excuse (a poor reason)\textquotedblright{}.

\begin{quote}
\emph{yàoshi nǐ bù xiǎng qù jiù bié zhǎo jièkǒu} — If you don't want to go, don't make excuses.
\end{quote}

Its characters come later; for now it is a word for the ear.

\begin{cousinweb}
\emph{jièkǒu} is built from \emph{jiè} 'to borrow' and \zh{口} 'a mouth', both inherited from Middle Chinese, \emph{tsjaeH} and \emph{khuwX} (Guangyun, in Baxter's transcription), which glosses the first 'to borrow': words borrowed to cover the real reason.
\end{cousinweb}

\begin{grammarlens}[title={What you've built}]
One more word for saying what would happen.
\end{grammarlens}

\subsection*{Your turn}
\begin{itemize}
\raggedright
  \item \emph{Say it:} \emph{jièkǒu}
  \item \emph{Say it:} \emph{jièkǒu}, once more
  \item \emph{Say it:} \emph{yàoshi nǐ bù xiǎng qù jiù bié zhǎo jièkǒu} — If you don't want to go, don't make excuses.
  \item \emph{Recall:} say the Mandarin for hard to come by, then the Mandarin for in an awkward spot, then say \emph{jièkǒu} again
\end{itemize}

\begin{tcolorbox}[breakable,colback=teal!4,colframe=teal!35!black,title={Before you move on}]
What does jièkǒu mean? (\textquotedblleft{}Excuse (a poor reason)\textquotedblright{}.) Say it once more.
\end{tcolorbox}
